



% \subsubsection{Tools}

% A \emph{tool} denotes an external function, API, or service that an LLM-based agent can invoke to perform actions beyond language generation. As summarized in Table~\ref{tab:framework-comparison}, tool support exists across all framework families but differs fundamentally in where execution authority resides and how invocation is controlled. In graph-based frameworks, tools are integrated procedurally either as agent-bound functions or as explicit nodes in the execution graph, with invocation permitted only at developer-specified locations and governed by static control flow~\cite{langgraph2025,langflow2025,autogpt2025,n8n2025}. This design emphasizes predictability and explicit execution ordering, but restricts agents from autonomously selecting or reconfiguring tools beyond the predefined workflow. Role-based frameworks associate tools with agent roles and task context, enabling dynamic invocation during task-driven reasoning or manager--worker coordination while still constraining access through role semantics and orchestration logic~\cite{crewai2025,wu2024autogen,openaiagents2025,OpenAgents}. In contrast, GABM frameworks do not expose tools to agents directly; instead, agents emit actions or intentions that are interpreted and executed by an environment controller, such as a Game Master, which applies tools to update global state~\cite{vezhnevets2023generative}. Across all frameworks, tool invocation relies on the LLM’s native function-calling interface, where tool signatures and parameters are provided to the model and selection is performed entirely by the LLM. No framework introduces explicit mechanisms for tool prefiltering or ranking, and all are constrained by model-level limits on the number of exposed tools (e.g., 128 tools in current OpenAI interfaces). As a result, tool integration functions as an execution facility rather than a first-class orchestration mechanism.





% \noindent\textbf{- Tool Execution Model.}
% This architectural dimension determines where authority for external action execution resides and how tool invocation is governed at runtime. Graph-based frameworks integrate tools procedurally in two ways. Tools may appear as explicit workflow nodes whose execution is fixed by control-flow logic~\cite{langgraph2025,langflow2025,n8n2025}. Alternatively, tools may be bound to agents as callable functions, but their invocation remains constrained by developer-defined graph structure~\cite{langgraph2025,autogpt2025}. Role-based frameworks bind tools to agent roles and task context, allowing dynamic selection during reasoning while constraining access through orchestration semantics~\cite{crewai2025,wu2024autogen,openaiagents2025,OpenAgents}. GABM frameworks remove direct tool access from agents entirely, routing all external actions through an environment controller that interprets agent intentions and updates global state~\cite{vezhnevets2023generative}. These designs trade execution predictability for adaptive autonomy. Procedural invocation offers strict control and stable cost, role-conditioned access enables flexible problem solving with weaker guarantees, and environment-mediated execution supports coherent long-horizon state evolution while limiting direct agent control. Accordingly, workflow-bound tools suit structured pipelines, role-scoped tools suit open-ended tasks requiring dynamic action choice, and environment-controlled tools suit simulation-driven systems with persistent global dynamics.

\smallskip
\noindent\textbf{- Tool Execution Model.}
This architectural dimension determines where the authority for external action execution resides and how tool invocation is governed at runtime. Graph-based frameworks integrate tools procedurally in two ways: as explicit workflow nodes with execution fixed by control-flow logic~\cite{langgraph2025,langflow2025,n8n2025}, or as agent-callable functions whose invocation remains constrained by developer-defined graph structure~\cite{langgraph2025,autogpt2025}. Role-based frameworks bind tools to agent roles and task context, enabling dynamic selection during reasoning while constraining access through orchestration semantics~\cite{crewai2025,wu2024autogen,openaiagents2025,OpenAgents}. In contrast, GABM frameworks remove direct tool access from agents and route all external actions through an environment controller that interprets agent intentions and updates global state~\cite{vezhnevets2023generative}. Together, these designs trade execution predictability for adaptive autonomy. Procedural invocation offers strict control and stable costs. Role-conditioned access enables flexible problem solving with weaker guarantees, while environment-mediated execution supports persistent state evolution while limiting direct agent control.\shorten

